% ============================================================
% SECTION 4: DATA ENGINEERING SOLUTION
% ============================================================
\section{Giải pháp Kỹ thuật Dữ liệu}
\label{sec:data-engineering-solution}

% --- 4.1 High-Level Architecture ---
\subsection{Kiến trúc tổng quan hệ thống}
\label{sec:architecture}

\begin{figure}[H]
    \centering
    \includegraphics[width=0.9\textwidth]{images/architecture.pdf}
    \caption[System Architecture]{Kiến trúc hệ thống tổng quan — Luồng dữ liệu từ tài xế đến giao diện khách hàng}
    \label{fig:system-architecture}
\end{figure}

\bigskip
\noindent
\textbf{Mô tả luồng dữ liệu (từ Hình \ref{fig:system-architecture}):}

\bigskip
\noindent
\textbf{Bước 1 — GPX Simulation (Golang):} Chương trình simulator đọc tệp GPX, trích xuất các điểm GPS, và gửi JSON event đến Kafka topic \texttt{raw-location-events}. Mỗi event được gắn key là \texttt{driver\_id} để đảm bảo tất cả event của cùng một tài xế được xử lý tuần tự trong cùng một partition.

\bigskip
\noindent
\textbf{Bước 2 — Stream Processing (Kafka Streams / Java):} Topology tiêu thụ \texttt{raw-location-events}, thực hiện: lọc tọa độ không hợp lệ, tính tốc độ qua State Store, tính tốc độ trung bình qua windowed aggregation (30s tumbling window), join với KTable điểm đích, phân nhánh alert. Kết quả: topic \texttt{processed-updates} và \texttt{alerts}.

\bigskip
\noindent
\textbf{Bước 3 — Real-Time Serving (Golang API + WebSocket):} API server (Gin) đồng thời tiêu thụ \texttt{processed-updates}, push qua WebSocket, và ghi \texttt{trip\_locations} vào Cassandra.

\bigskip
\noindent
\textbf{Bước 4 — Frontend (React + Leaflet):} Ứng dụng web hiển thị bản đồ OpenStreetMap, cập nhật vị trí tài xế mỗi giây qua WebSocket.

% --- 4.2 Ingestion Layer ---
\subsection{Tầng nhập liệu (Golang)}
\label{sec:ingestion-layer}

\begin{figure}[H]
    \centering
    \includegraphics[width=0.7\textwidth]{images/golang-ingestion.pdf}
    \caption[Ingestion Layer Architecture]{Kiến trúc tầng nhập liệu — Goroutine cho mỗi tài xế}
    \label{fig:ingestion-architecture}
\end{figure}

\bigskip
\noindent
Mỗi goroutine xử lý một tài xế: \texttt{gpx.Parse(file)} đọc toàn bộ GPX → vòng lặp duyệt từng điểm GPS, tạo \texttt{LocationEvent}, gọi \texttt{producer.Publish()} → \texttt{time.Sleep(interval)} chờ giữa 2 điểm (mặc định 1000ms).

\bigskip
\noindent
\textbf{Tại sao Golang?}
\begin{itemize}
    \item \textbf{Goroutine lightweight:} ~2KB stack mỗi goroutine — một laptop có thể chạy 10.000+ goroutines đồng thời.
    \item \textbf{Type safety:} Struct \texttt{LocationEvent} định nghĩa tại compile-time, đảm bảo JSON đúng schema.
    \item \textbf{Performance:} Compiled language — parse XML và serialize JSON nhanh.
\end{itemize}

% --- 4.3 Processing Layer (Kafka Streams) ---
\subsection{Tầng xử lý luồng (Kafka Streams — Java)}
\label{sec:processing-layer}

Kafka Streams là thư viện xử lý luồng chạy trong JVM. Phần này trình bày chi tiết topology được triển khai.

\begin{enumerate}
    \item \textbf{Filter Invalid Coordinates (RT1):} Loại bỏ tọa độ lat/lng nằm ngoài phạm vi, accuracy > 20m.
    \item \textbf{Stateful Speed Calculation (RT3):} \texttt{SpeedAlertProcessor} sử dụng State Store (RocksDB) lưu GPS point trước đó. Tính \texttt{velocity = (dist\_km / time\_delta\_s) × 3600}. Nếu > 60 km/h → cờ \texttt{isSpeeding = true}.
    \item \textbf{Windowed Aggregation (RT2):} \texttt{TimeWindows.ofSizeWithNoGrace(Duration.ofSeconds(30))} — cửa sổ 30 giây không overlap. Tính tốc độ trung bình.
    \item \textbf{Stream-Table Join (RT4):} \texttt{KTable<String, Destination>} từ topic \texttt{orders}. Join mỗi location event với destination tương ứng → tính Haversine distance.
    \item \textbf{Branching:} Chia luồng: SPEEDING alert / PROXIMITY alert / valid location.
\end{enumerate}

\begin{lstlisting}[language=Java, caption={Kafka Streams Topology — \texttt{Main.java}: Filter $\rightarrow$ Stateful Speed $\rightarrow$ Windowed Avg $\rightarrow$ Stream-Table Join $\rightarrow$ Branch}, label={lst:topology}]
public static Topology buildTopology() {
    StreamsBuilder builder = new StreamsBuilder();

    // Serde definitions (JSON)
    JsonSerde<LocationEvent>       locationEventSerde = new JsonSerde<>(LocationEvent.class);
    JsonSerde<EnrichedLocation>    enrichedLocationSerde = new JsonSerde<>(EnrichedLocation.class);
    JsonSerde<Alert>              alertSerde            = new JsonSerde<>(Alert.class);
    JsonSerde<Destination>         destinationSerde       = new JsonSerde<>(Destination.class);
    JsonSerde<SpeedAccumulator>   speedAccumulatorSerde  = new JsonSerde<>(SpeedAccumulator.class);

    // 1. Add state store for speed calculation (RocksDB-backed)
    builder.addStateStore(Stores.keyValueStoreBuilder(
        Stores.persistentKeyValueStore(SpeedAlertProcessor.STATE_STORE),
        Serdes.String(), locationEventSerde));

    // Consume raw events, re-key by driver_id for partitioning
    KStream<String, LocationEvent> rawStream = builder.stream(
            "raw-location-events",
            Consumed.with(Serdes.String(), locationEventSerde))
        .selectKey((k, v) -> v != null ? v.getDriverId() : "unknown");

    // 2. Filter Invalid Coordinates (RT1)
    KStream<String, LocationEvent> filteredStream = rawStream.filter(
        (key, event) ->
            event != null &&
            event.getLatitude()  >= -90 && event.getLatitude()  <= 90 &&
            event.getLongitude() >= -180 && event.getLongitude() <= 180 &&
            event.getAccuracy()  <= 20.0);

    // 3. Stateful Speed Calculation (RT3) via Processor API
    KStream<String, EnrichedLocation> speedStream = filteredStream.process(
        new SpeedAlertProcessor(), SpeedAlertProcessor.STATE_STORE);

    // 4. Windowed Aggregation -- 30s tumbling window (RT2)
    KTable<Windowed<String>, Double> avgSpeedTable = speedStream
        .groupByKey(Grouped.with(Serdes.String(), enrichedLocationSerde))
        .windowedBy(TimeWindows.ofSizeWithNoGrace(Duration.ofSeconds(30)))
        .aggregate(
            SpeedAccumulator::new,
            (key, value, agg) -> agg.add(value.getSpeed()),
            Materialized.<String, SpeedAccumulator,
                 org.apache.kafka.streams.state.WindowStore
                 <org.apache.kafka.common.utils.Bytes,byte[]>>
                 as("speed-window-store")
                 .withValueSerde(speedAccumulatorSerde))
        .mapValues(SpeedAccumulator::getAverage);

    // 5. Stream-Table Join with destination (RT4)
    KTable<String, Destination> destinationTable = builder.table(
        "orders",
        Consumed.with(Serdes.String(), destinationSerde),
        Materialized.as("destination-store"));

    KStream<String, EnrichedLocation> enrichedStream = speedStream.leftJoin(
        destinationTable,
        (loc, dest) -> {
            if (dest != null) {
                double distKm = Haversine.haversine(
                    loc.getLatitude(), loc.getLongitude(),
                    dest.getLatitude(), dest.getLongitude());
                int etaSec = (loc.getSpeed() > 0)
                    ? (int) ((distKm / loc.getSpeed()) * 3600) : 0;
                return loc.withDistanceAndEta(distKm, etaSec);
            }
            return loc;
        },
        Joined.with(Serdes.String(), enrichedLocationSerde, destinationSerde));

    // 6. Branch: Speeding | Proximity (<500m) | Valid
    KStream<String, EnrichedLocation>[] branches = enrichedStream.branch(
        (k, v) -> v.isSpeeding(),                             // Branch 0: SPEEDING
        (k, v) -> v.getDistanceToDestination() > 0
               && v.getDistanceToDestination() < 0.5,          // Branch 1: PROXIMITY
        (k, v) -> true);                                      // Branch 2: VALID

    // Branch 0 -\textgreater{} alerts (SPEEDING)
    branches[0].map((k, loc) -> new KeyValue<>(k,
        new Alert(UUID.randomUUID().toString(), loc.getDriverId(),
                 loc.getTripId(), Instant.now(), "SPEEDING", "HIGH",
                 String.format("Speed limit exceeded: %.0f km/h in 60 km/h zone",
                               loc.getSpeed()),
                 Map.of("current_speed", String.valueOf((int) loc.getSpeed()),
                        "limit", "60",
                        "location", loc.getLatitude()+","+loc.getLongitude()))))
        .to("alerts", Produced.with(Serdes.String(), alertSerde));

    // Branch 1 -\textgreater{} alerts (PROXIMITY)
    branches[1].map((k, loc) -> new KeyValue<>(k,
        new Alert(UUID.randomUUID().toString(), loc.getDriverId(),
                 loc.getTripId(), Instant.now(), "PROXIMITY", "MEDIUM",
                 String.format("Driver is %.0fm away",
                               loc.getDistanceToDestination() * 1000),
                 Map.of("distance_m",  String.valueOf((int)(loc.getDistanceToDestination()*1000)),
                        "threshold_m","500"))))
        .to("alerts", Produced.with(Serdes.String(), alertSerde));

    // Branch 2 -\textgreater{} processed-updates
    branches[2].to("processed-updates",
        Produced.with(Serdes.String(), enrichedLocationSerde));

    return builder.build();
}
\end{lstlisting}

\begin{lstlisting}[language=Java, caption={SpeedAlertProcessor — Stateful Processor với State Store (RocksDB)}, label={lst:speed-processor}]
public class SpeedAlertProcessor
        implements Processor<String, LocationEvent, String, EnrichedLocation> {

    public static final String STATE_STORE = "previous-location-store";
    private ProcessorContext context;
    private StateStore stateStore;

    @Override
    public void init(ProcessorContext context) {
        this.context  = context;
        this.stateStore = context.getStateStore(STATE_STORE);
    }

    @Override
    public void process(String key, LocationEvent event) {
        // Read previous GPS point from RocksDB state store
        LocationEvent prev = (LocationEvent) stateStore.get(key);

        if (prev != null) {
            double distKm = Haversine.haversine(
                    prev.getLatitude(),  prev.getLongitude(),
                    event.getLatitude(), event.getLongitude());

            double timeDeltaSec = (event.getTimestamp().toEpochMilli()
                                 - prev.getTimestamp().toEpochMilli()) / 1000.0;

            if (timeDeltaSec > 0) {
                // velocity = (km / seconds) * 3600 = km/h
                double speedKmh = (distKm / timeDeltaSec) * 3600;
                context.forward(key, new EnrichedLocation(event, speedKmh));
            }
        }

        // Update state store for next computation
        stateStore.put(key, event);
    }

    @Override public void close() { }
}
\end{lstlisting}

% --- 4.4 Storage Layer (Cassandra) ---
\subsection{Tầng lưu trữ (Apache Cassandra)}
\label{sec:storage-layer}

Apache Cassandra được chọn cho khả năng ghi cực cao và mở rộng ngang tuyến tính. Phần này trình bày chi tiết thiết kế schema CQL.

\begin{lstlisting}[language=SQL, caption={Cassandra Schema — \texttt{init-cql.cql}: Keyspace và 5 bảng chính}, label={lst:cql-schema}]
-- ============================================================
-- Keyspace
-- ============================================================
CREATE KEYSPACE IF NOT EXISTS delivery_tracking
WITH REPLICATION = {'class': 'SimpleStrategy', 'replication_factor': 1};

USE delivery_tracking;

-- ============================================================
-- Table 1: orders
-- Query: SELECT * FROM orders WHERE order_id = ?
-- ============================================================
CREATE TABLE IF NOT EXISTS orders (
    order_id             UUID PRIMARY KEY,
    customer_id         UUID,
    driver_id           TEXT,
    restaurant_location TEXT,
    delivery_location   TEXT,
    status              TEXT,
    created_at          TIMESTAMP,
    updated_at          TIMESTAMP
);

-- ============================================================
-- Table 2: trip_locations  (TIME-SERIES)
-- Query: SELECT * FROM trip_locations
--        WHERE trip_id = ? ORDER BY timestamp ASC
-- ============================================================
CREATE TABLE IF NOT EXISTS trip_locations (
    trip_id    UUID,
    timestamp  TIMESTAMP,
    driver_id  TEXT,
    order_id   UUID,
    latitude   DOUBLE,
    longitude  DOUBLE,
    speed      DOUBLE,
    heading    DOUBLE,
    accuracy   DOUBLE,
    PRIMARY KEY (trip_id, timestamp)
) WITH CLUSTERING ORDER BY (timestamp DESC);

-- ============================================================
-- Table 3: trip_metadata  (AGGREGATE)
-- Query: SELECT * FROM trip_metadata WHERE trip_id = ?
-- ============================================================
CREATE TABLE IF NOT EXISTS trip_metadata (
    trip_id             UUID PRIMARY KEY,
    driver_id           TEXT,
    order_id            UUID,
    start_time          TIMESTAMP,
    end_time            TIMESTAMP,
    start_location      TEXT,
    destination         TEXT,
    total_distance      DOUBLE,
    total_duration      INT,
    average_speed       DOUBLE,
    max_speed           DOUBLE,
    speeding_violations INT,
    trip_cost           DECIMAL,
    status              TEXT
);

-- ============================================================
-- Table 4: driver_analytics  (WEEKLY AGGREGATION)
-- Query: SELECT * FROM driver_analytics
--        WHERE driver_id = ? AND week_start_date >= '2024-01-01'
-- ============================================================
CREATE TABLE IF NOT EXISTS driver_analytics (
    driver_id           TEXT,
    week_start_date     DATE,
    total_trips         INT,
    total_distance      DOUBLE,
    total_duration      INT,
    average_speed       DOUBLE,
    speeding_violations INT,
    idle_time           INT,
    PRIMARY KEY (driver_id, week_start_date)
) WITH CLUSTERING ORDER BY (week_start_date DESC);

-- ============================================================
-- Table 5: alerts  (AUDIT TRAIL)
-- Query: SELECT * FROM alerts
--        WHERE driver_id = ? AND timestamp >= ?
-- ============================================================
CREATE TABLE IF NOT EXISTS alerts (
    alert_id    UUID,
    driver_id   TEXT,
    trip_id     UUID,
    timestamp   TIMESTAMP,
    alert_type  TEXT,
    severity    TEXT,
    message     TEXT,
    metadata    MAP<TEXT, TEXT>,
    PRIMARY KEY (driver_id, timestamp, alert_id)
) WITH CLUSTERING ORDER BY (timestamp DESC);

CREATE INDEX IF NOT EXISTS ON trip_metadata (order_id);
\end{lstlisting}

\bigskip
\noindent
\textbf{Lý do chọn PRIMARY KEY:}

\bigskip
\begin{itemize}
    \item \texttt{trip\_locations — PRIMARY KEY (trip\_id, timestamp):}
    \begin{itemize}
        \item \textbf{Partition Key = \texttt{trip\_id}:} Tất cả điểm GPS của cùng một chuyến đi nằm trên cùng một Cassandra node. Truy vấn \texttt{WHERE trip\_id = ?} chỉ cần đọc từ 1 node.
        \item \textbf{Clustering Key = \texttt{timestamp DESC}:} Điểm mới nhất ở đầu partition. Truy vấn \texttt{LIMIT 1} trả về vị trí mới nhất ngay lập tức.
    \end{itemize}

    \item \texttt{driver\_analytics — PRIMARY KEY (driver\_id, week\_start\_date):}
    \begin{itemize}
        \item \textbf{Partition Key = \texttt{driver\_id}:} Toàn bộ dữ liệu của một tài xế nằm trên 1 partition.
        \item \textbf{Clustering Key = \texttt{week\_start\_date DESC}:} Tuần mới nhất trước — phù hợp pattern "xem tuần gần đây".
    \end{itemize}

    \item \texttt{alerts — PRIMARY KEY (driver\_id, timestamp, alert\_id):}
    \begin{itemize}
        \item \textbf{Partition Key = \texttt{driver\_id}:} Mỗi driver có 1 partition riêng.
        \item \textbf{Clustering Key = \texttt{timestamp DESC}:} Alert mới nhất luôn ở đầu. \texttt{alert\_id} đảm bảo uniqueness khi 2 alert cùng timestamp.
    \end{itemize}
\end{itemize}
